	1. Introducción breve del hallazgo

Contextualiza qué parte del objetivo aborda el hallazgo (ej. frecuencia de términos, tópicos emergentes, valoración discursiva).

Usa un enunciado claro, tipo:
“En esta sección presentamos hallazgos preliminares sobre la frecuencia y los contextos en los que aparece el término ciencia en las columnas de opinión analizadas (2018–2020).”

\begin{itemize}
	\item \textbf{Coocurrencias y redes semánticas:} PMI/PPMI; grafos conceptos–actores.
	\item \textbf{Tópicos:} BERTopic/NMF con validación por coherencia y revisión humana.
	\item \textbf{Valoración/stance:} esquema de anotación (acuerdo interanotador), clasificador ligero (LogReg/SVM) o modelo BETO/roberta-es; métricas (F1, AUC).
	\item \textbf{Dinámica temporal:} series mensuales y detección de puntos de cambio.
\end{itemize}

2. Evidencia cuantitativa

Gráficos o tablas simples:

Tendencias temporales (ej. apariciones de ciencia por mes/año).

Distribución por medio o autor.

Top 20 términos que coocurren con ciencia.

Ejemplo de redacción:
“Entre 2018 y 2020, el término ciencia aparece en 2.340 columnas (29,2% del corpus). Su presencia es más recurrente en 2020, coincidiendo con la pandemia de COVID-19.”

3. Evidencia cualitativa

Fragmentos de columnas donde la palabra aparece con funciones diferentes (argumentativa, retórica, simbólica).

Ejemplo:
“En algunos casos, ciencia se emplea como fuente de legitimidad (‘la ciencia demuestra que…’), mientras que en otros se usa con valor metafórico (‘la ciencia política de la corrupción’).”

4. Conexión con la pregunta de investigación

Relaciona el hallazgo con las preguntas planteadas en el resumen:

¿Qué usos del discurso científico identificamos hasta ahora?

¿Cómo se valora la ciencia en la esfera pública?

Ejemplo:
“Estos hallazgos preliminares sugieren que la ciencia no solo se asocia a temas de salud y educación, sino que también funciona como metáfora en discusiones políticas, lo cual abre la pregunta sobre el grado en que la ciencia es apropiada simbólicamente en la opinión pública.”

5. Visualización / síntesis

Incluye una gráfica o tabla por hallazgo preliminar:

Nube de palabras coocurrentes con ciencia.

Gráfico de barras con los tópicos más frecuentes.

Red semántica pequeña con actores y conceptos conectados.

Ejemplo de redacción de un hallazgo preliminar:

Hallazgo preliminar 1: Aparición y contextos del término ciencia

Del total de 8.000 columnas analizadas, 2.340 (29,2%) incluyen al menos una mención explícita a ciencia. La frecuencia muestra un aumento pronunciado en 2020, en paralelo con la pandemia.

En términos de coocurrencias, ciencia aparece asociada principalmente a salud, educación y tecnología. Sin embargo, también se emplea en contextos metafóricos vinculados a la política.

Estos resultados preliminares sugieren que la ciencia cumple un doble papel en la esfera pública: como autoridad epistémica (ej. validación de argumentos) y como recurso retórico en debates no especializados.
